\scanpage{85}
Et en particulier, si $g$ de la forme $\sigma_a^0$, on trouve
\[
(82)\qquad
\left\{
\begin{aligned}
\widetilde\sigma_a(\varepsilon_s)&=\varepsilon_{\sigma_a^0(s)}^{-1}\\
\widetilde\sigma_a(\rho_s)&=\rho_{\sigma_a^0(s)}^{-1}
\end{aligned}
\right.
\qquad s\in S,\ a\in A
\]
\margin{donc $\widetilde\sigma_{a_0}(\rho_i)=\rho_{1-i}^{-1}$}
ce qui donne l'opération des $\widetilde\sigma_a$ sur $\Pi_P$, \deleted{\unsure{bon}} d'où \unsure{en principe} \deleted{\unsure{leur}} (avec celle des $\tau_a$) l'opération du $\Pi_P^{D,\tau}$ entier sur $\Pi_P$. \inserted{comme $\sigma_{a_0}=\widetilde\sigma_{a_0}\tau_{a_0}$} En particulier,
\[
\left\{
\begin{aligned}
\sigma_{a_0}(\rho_0)&=\rho_1,\quad
\sigma_{a_0}(\rho_1)=\rho_0,\quad
\sigma_{a_0}(\rho_i)=\tau_{a_0}(\rho_{n+1-i}^{-1})\\
&\hspace{13em}=(\rho_n\cdots\rho_{n+2-i})(\rho_{n+1-i})\\
&\hspace{13em}\text{pour }2\leq i\leq n-1.
\end{aligned}
\right.
\]
\margin{\uline{NB} L'action de $\sigma_a$ sur $\Pi_P$ est la composée de
\[
\Pi_1(\sigma_a^0):\Pi_1(\Sigma^*,P_\omega)\longrightarrow\Pi_1(\Sigma^*,P_{\omega'})
\]
avec
\[
\lambda_a:\Pi_1(\Sigma^*,P_{\omega'})\longrightarrow\Pi_1(\Sigma^*,P_\omega).
\]}

\deleted{\unsure{On}} On trouve d'ailleurs \deleted{$\Pi_P^\tau$}
\[
(\unsure{84|83})\qquad
\Pi_{P_{\unsure{c}}}^{D,\tau}\simeq\Pi_P^D\cdot\{1,\tau_{a_0}\}
\qquad\text{produit }\tfrac12\text{ direct},
\]
il faut expliciter encore l'opération de $\tau_{a_0}$ sur $\Pi_P^D$ i.e. sur $\rho$ et sur $\sigma_{a_0}$, on a
\[
(85)\qquad
\tau_{a_0}(\rho)=\sigma_{a_0}(\rho^{-1}),\quad
\tau_{a_0}(\sigma_{a_0})=\sigma_{a_0},
\]
la deuxième relation est déjà notée (72), la première s'obtient en écrivant
\[
\deleted{\rho}\tau_{a_0}
\underbrace{\rho\tau_{a_0}\rho^{-1}}_{\tau_{a_1}}
=\tau_{a_0}\tau_{a_1}
=\rho_1^{-1}
=(\rho\sigma_{a_0}\rho\sigma_{a_0})^{-1}
=\sigma_{a_0}\rho^{-1}\sigma_{a_0}\rho^{-1}
\]
d'où (85) ; \deleted{\unsure{qui s'écrit aussi}}\inserted{plus simplement} elle s'écrit
\[
\widetilde\sigma_{a_0}(\rho)=\rho^{-1}
\]
qui est triviale, puisque \deleted{pour} $\rho=\widetilde\omega$ \deleted{\unsure{de}} $g\mapsto\widetilde g$ est un hom de groupes $D_P\mathbin{-}\Pi_1^{P,\tau}$, enfin $\sigma_{a_0}^0(\rho)=\rho^{-1}\ldots$.
Il est plus joli peut-être d'écrire
\[
\textcircled{86}\quad
\left\{
\begin{aligned}
\Pi_P^{D,\tau}&\simeq\Pi_P^D\cdot\{1,\widetilde\sigma_{a_0}\}\\
\widetilde\sigma_{a_0}(\rho)&=\rho^{-1},\quad
\widetilde\sigma_{a_0}(\sigma_{a_0})=\sigma_{a_0}
\end{aligned}
\right.
\quad
\left[
\begin{gathered}
=\{\rho,\sigma_{a_0},\widetilde\sigma_{a_0}\mid
\rho^n=\sigma_{a_0}^2=\widetilde\sigma_{a_0}^{\,2}
=(\sigma_{a_0}\widetilde\sigma_{a_0})^2\\
=\widetilde\sigma_{a_0}(\rho)\rho=1\}
\end{gathered}
\right]
\]
\margin{84}
